\section{Collective curvature and a graph-resolved finite remainder}
\label{sec:walkbounds}
\status{EXACT\_IDENTITY}
Let $d\in\mathbb R^{3n}$ denote ordinary replacement and gain increments.
Define fixed diagonal selectors $E_a$ such that
$\Theta(d)=\sum_a d_a E_a$: a replacement selector contains two unit
entries, whereas a gain selector contains one. In the affine descriptor
chart, $H(s,d)=\sum_a d_a E_aT(s)$ exactly.
All delays remain fixed inside $T(s)$.

\begin{theorem}[Collective cluster curvature]
Under the contour hypotheses of Section~\ref{sec:regional}, the real
cluster margin has the second-order expansion
\begin{equation}
 g_\Gamma(p_0+d)=g_{\Gamma,0}+J_\Gamma d+
       \tfrac12d^TQ_\Gamma d+\mathcal R_{\Gamma,3}(d),\qquad
 (Q_\Gamma)_{ab}=
 \Re\frac{1}{2\pi\ii m_\Gamma}
       \oint_\Gamma\operatorname{tr}(E_aTE_bT)\,ds .
 \label{eq:collective-hessian}
\end{equation}
The matrix $Q_\Gamma$ is real symmetric. It is not required to be
positive semidefinite and is not a Laplacian.
\end{theorem}
\begin{proof}
Substitute $H=\sum_a d_aE_aT$ in the convergent series for
$-\operatorname{tr}\log(I+H)=-\operatorname{tr}H+
\operatorname{tr}H^2/2-\cdots$ in \eqref{eq:regional-log}.
Cyclic invariance of the trace gives
$\operatorname{tr}(E_aTE_bT)=\operatorname{tr}(E_bTE_aT)$.
Collecting terms proves the formula. Uniform convergence on the
compact contour permits termwise integration. This differentiates a
cluster moment, so internal root collisions do not require eigenvector
derivatives or a simple-root assumption.
\end{proof}
For scalar selectors the integrand is $T_{ab}T_{ba}$.
Thus a mixed change at two sites follows a return path through the
full dynamical network. Two individually favorable first-order
changes need not remain favorable when combined. The sign is set by
the complex contour integral, not by graph adjacency or a positive
edge weight. Diagonal terms already include the full grid.

\status{PROVED\_REDUCED\_MODEL}
Suppose a continuous, piecewise bounded nonnegative envelope satisfies
\[
 |H(s,d)|\le M(s),\qquad
 \sup_{s\in\Gamma}\rho(M(s))<1,\qquad d\in\mathcal D_d .
\]
A verified positive-vector inequality on every panel is one sufficient
way to establish this condition. Set
\[
 \Psi_k(M)=\sum_{\ell=k}^{\infty}\frac{\operatorname{tr}M^\ell}{\ell}
 =-\log\det(I-M)-
     \sum_{\ell=1}^{k-1}\frac{\operatorname{tr}M^\ell}{\ell},
 \qquad k\ge2 .
\]
Here the determinant is positive and the logarithm is real: real
eigenvalues of $M$ are below one, and nonreal eigenvalues occur in
conjugate pairs. Despite that spectral representation, every term in
the defining trace series is nonnegative.

\begin{theorem}[Closed-walk remainder and interaction-omission budget]
The remainder after degree $k-1$ obeys
\begin{equation}
 \boxed{|\mathcal R_{\Gamma,k}(d)|\le
 \beta_{\Gamma,k}^{\rm walk}:=
 \frac{1}{2\pi m_\Gamma}
       \oint_\Gamma\Psi_k(M(s))\,|ds|.}
 \label{eq:walk-remainder}
\end{equation}
Partition channels into physical intervention sites and let
$H_{\rm b}=\operatorname{blockdiag}(H_{11},\ldots,H_{nn})$.
The omitted contribution to the cluster-shift formula, obtained by
replacing $H$ with $H_{\rm b}$, is bounded by
\begin{equation}
 \boxed{\beta_\Gamma^{\rm cross}=
 \frac{1}{2\pi m_\Gamma}\oint_\Gamma
 \left[\Psi_2(M)-\sum_i\Psi_2(M_{ii})\right]|ds|.}
 \label{eq:cross-budget}
\end{equation}
This is a bound on an algebraic approximation to the contour formula;
it does not assert that $H_{\rm b}$ is a realizable physical grid.
\end{theorem}
\begin{proof}
The entrywise inequality $|H^\ell|\le M^\ell$ implies
$|\operatorname{tr}H^\ell|\le\operatorname{tr}M^\ell$.
Sum the tail of the logarithm series and integrate its absolute value
to obtain \eqref{eq:walk-remainder}.
For \eqref{eq:cross-budget}, the terms remaining after block-diagonal
subtraction are exactly index sequences that visit more than one
block before returning to their starting channel. Bounding the
absolute value of each product by the corresponding product of
entries of $M$ gives
$\operatorname{tr}M^\ell-\sum_i\operatorname{tr}M_{ii}^\ell$.
The degree-one difference is zero. Summation proves the result.
\end{proof}

The earlier scalar bound follows from
$\operatorname{tr}M^\ell\le r\kappa^\ell$ whenever a common induced
norm of $M$ is bounded by $\kappa<1$. Retaining the actual walk sums
can therefore substantially reduce conservatism. A computable,
slightly looser expression avoids logarithms:
\[
 \Psi_k(M)\le\frac1k\operatorname{tr}\!\left[M^k(I-M)^{-1}\right].
\]
All inverse, determinant, integration and rounding errors must still
be enclosed. In particular, a verified evaluation of
\eqref{eq:cross-budget} requires an \emph{upper} enclosure of the
total minus \emph{lower} enclosures of the within-site terms;
subtracting two upper bounds would be invalid. Directly summing
cross-site walks is an alternative.

\paragraph{Meaning of beyond nodal damping.}
Equation~\eqref{eq:collective-hessian} identifies the signed,
frequency-dependent interaction between two replacement or tuning
actions. Equation~\eqref{eq:cross-budget} quantifies the error made
by dropping those interactions. These are not nodal damping
constants: $T$ contains the grid resolvent, SG dynamics, voltage/PQ
coupling and the exact phases $e^{-s\tau_i}$. They supplement the
physical torque--speed damping operator derived in the appendices;
they do not turn a contour return matrix into a physical impedance.
The closed walks are paths in the intervention-port representation,
not automatically cycles of transmission lines. No communication
between PLLs has been introduced.

\subsection{A stronger necessary replacement relaxation}
Write $d\in[l,u]$ and introduce a real symmetric matrix $Z$.
Every actual design admits $Z=dd^T$, hence satisfies
\begin{align}
 &\begin{bmatrix}1&d^T\\d&Z\end{bmatrix}\succeq0,\\
 &Z_{ab}\ge l_ad_b+l_bd_a-l_al_b,\qquad
 Z_{ab}\ge u_ad_b+u_bd_a-u_au_b,\\
 &Z_{ab}\le u_ad_b+l_bd_a-u_al_b,\qquad
 Z_{ab}\le l_ad_b+u_bd_a-l_au_b .
\end{align}
Retain the diagonal case of these inequalities as well.
Impose, for each verified cluster,
\begin{equation}
 g_{\Gamma,0}+J_\Gamma d+
       \tfrac12\langle Q_\Gamma,Z\rangle
       \le\beta_{\Gamma,3}^{\rm walk}.
 \label{eq:quadratic-outer}
\end{equation}
Maximizing $(P^0)^T\delta\rho$ over this semidefinite relaxation gives
an optimistic regional replacement upper bound. Dropping
$Z=dd^T$ enlarges the feasible set, which is the required direction;
the returned gains need not be physically feasible.
Neither a primal floating-point optimum nor a tiny solver residual
is a certified upper bound. A feasible dual with enclosed errors,
or another verified objective bound, is required for that label.

Certified coefficient errors add
$e_0+\sum_a e_{J,a}r_a+
\tfrac12\sum_{a,b}e_{Q,ab}r_ar_b$ to the right-hand side,
where $r_a=\max(|l_a|,|u_a|)$. The lift may sharpen the linear
relaxation, but no universal dominance follows if different
remainders or additional relaxations are used. Both outer problems
can be intersected because each contains all secure designs.
This strengthening has been derived and identity-tested; the
IEEE--39 semidefinite upper-bound problem has not been solved.
